The application objective is to restore the coupled nominal magnetic
motion through learned residual and damping compensation. The physical
input map must realize the feedback in
\eqref{struct:route-a-control}, preserve its estimation errors and use
the same midpoint convention as \eqref{ctrl:midpoint}. We first close
that chain for the two-mode object, then expose the additional terms
when the same fitted controller acts on more beam modes. The
full-beam clearance argument is a separate condition on physical
energy. Tracking a learned reference is retained as an extension in
Supplementary Section~\ref{sec:supp-tracking}.

\subsection{Magnetic realization of the learned residual compensation}
\label{app:ports}
Use $U=U_A=\widehat\gamma Y+\nabla\widehat R(X)$ from
\eqref{struct:route-a-control} and solve
\begin{equation}
 \binom{F}{\tau_B}=A_{\rm mag}^{-1}\Lambda U.
 \label{app:port-inverse}
\end{equation}
The continuous input $U\in\mathbb R^2$ is dimensionless; $F$ and
$\tau_B$ have units N and N\,m. The four-dimensional discrete increment
in \eqref{ctrl:midpoint} also contains the time step. These conversions
keep the physical force, the continuous generalized input and the
phase-space increment distinct.

Here $\tau_B$ is the active magnetic torque. In the two-mode physical
variables, the synthetic plant used in Section~\ref{sec:analysis} is
\begin{align}
 \ddot a+\alpha\dot a+\Omega_2^2a
 &=\boldsymbol t E_0(\vartheta-\sin\vartheta)
       -\Lambda\nabla R_*(X)\nonumber\\
 &\quad+A_{\rm mag}\binom{F}{\tau_B},
 \qquad a=D_sX.
 \label{app:residual-plant}
\end{align}
Since $\ddot a=\Lambda Y'$ and
$\Lambda\nabla V_0=\Omega_2^2a+\boldsymbol tE_0(\sin\vartheta-\vartheta)$,
substitution recovers \eqref{struct:plant}. A total modal moment is an
equivalent port convention:
\begin{align}
 T&=E_0(\vartheta-\sin\vartheta)+\tau_B,\nonumber\\
 A_{\rm mag}\binom{F}{T}
 &=\Lambda U+\boldsymbol tE_0(\vartheta-\sin\vartheta).
 \label{app:total-port}
\end{align}
The bias-field remainder is known, while the true residual force
$-\Lambda\nabla R_*$ remains in the plant. The controller does not
read $R_*$ through this change of convention.

Substituting the learned law into the physical plant now gives
\begin{align}
 \ddot a&=-\Lambda\nabla V_0
       +(\widehat\alpha-\alpha)\dot a\nonumber\\
       &\quad-\Lambda\nabla\delta R,\nonumber\\
 \widehat\alpha&=\omega_0\widehat\gamma,\qquad
 \delta R=R_*-\widehat R.
 \label{app:physical-closed-loop}
\end{align}
Thus the physical closed loop and
\eqref{struct:route-a-closed} carry the same potential and damping
errors. The compensated target retains the complete coupled nominal
magnetic potential. The truth-only comparator uses known synthetic
parameters solely to separate numerical error from identification
error.

Let the planar dipole be
$m=\mu(\cos\vartheta,\sin\vartheta,0)^{\mathsf T}$.
A local field with the required force and torque is
\begin{align}
 B(r)&=(B_0,B_y,0)^{\mathsf T}+G(r-r_{\rm tip}),\nonumber\\
 B_y&=\frac{\tau_B}{\mu\cos\vartheta},\nonumber\\
 G&=\frac{F}{\mu}
 \begin{pmatrix}
 -\sin\vartheta&\cos\vartheta&0\\
 \cos\vartheta&\sin\vartheta&0\\
 0&0&0
 \end{pmatrix}.
 \label{app:field-gradient}
\end{align}
Direct multiplication gives $Gm=(0,F,0)^{\mathsf T}$ and
$\|G\|_2=|F|/\mu$. Symmetry and zero trace satisfy the local
source-free curl and divergence conditions \citep{abbott}.
The active torque is $\mu B_y\cos\vartheta$; the axial bias torque
$-\mu B_0\sin\vartheta$ remains in the stored-energy model.

Every compensation ablation passes through this same interface.
For the learned group, reconstruct the applied transverse force and active torque from the
field and gradient, and form
\begin{align}
 U_{\rm rec}&=\Lambda^{-1}A_{\rm mag}
                    \binom{F_{\rm app}}{\tau_{B,\rm app}},\nonumber\\
 e_{\rm port}&=U_{\rm rec}-U_A(X,Y).
 \label{app:input-reconstruction}
\end{align}
Ideal exact conversion gives $e_{\rm port}=0$. Calibration errors,
angle-estimation errors, current limits and saturation can make this
quantity nonzero; an input-reconstruction check must retain them in
the applied plant. The reported tests verify the local algebraic
conversion and do not implement saturation.

The interface requires $A_{\rm mag}$ to have rank two, an angle domain
with $|\cos\vartheta|$ bounded away from zero, and feasible joint
field--gradient commands. Column normalization of the nominal matrix
gives condition number $3.8186$, a geometric diagnostic. A physical
authority calculation instead scales the columns by available force
and torque and uses an identified coil-current allocation. A single
uniform field supplies only one modal input direction.

Bias and transverse fields combine as vector magnitudes: a $30$ mT
bias with a $6$ mT transverse component requires $30.594$ mT, whereas
a $20$ mT transverse component requires $36.056$ mT. The published
$40$ mT arbitrary-direction capability of the OctoMag apparatus
\citep{lussi2021cvs} supports this field scale without determining
simultaneous gradient, workspace or bandwidth limits for the present
design. The learned residual-compensation test requires a peak
transverse field of approximately $0.54449$ mT. Section~\ref{sec:analysis}
reports its input reconstruction and sampled operating-domain checks.

\subsection{The same controller on an enriched modal plant}
\label{app:full-model}
Fix the mass-normalized eigenbasis of the $32$-element beam. Retaining
$N\in\{2,4,8,64\}$ modes changes the plant while the fitted two-mode
feedback stays fixed. Write
$a^{(N)}=(a_1,\ldots,a_N)^{\mathsf T}$, and let
$\boldsymbol b_N,\boldsymbol t_N$ collect the corresponding tip
coefficients. Define
\begin{align}
 \vartheta_N&=\boldsymbol t_N^{\mathsf T}a^{(N)},\nonumber\\
 A_{\rm mag}^{(N)}&=[\boldsymbol b_N\ \boldsymbol t_N]
                         \in\mathbb R^{N\times2},\nonumber\\
 X&=D_s^{-1}a_{1:2},\nonumber\\
 Y&=\omega_0^{-1}\mathsf N^{-1}D_s^{-1}\dot a_{1:2}.
 \label{app:enriched-coordinates}
\end{align}
Here $D_s,\mathsf N$ and $A_{\rm mag}$ still denote the fixed low-mode
maps. The first two rows of $A_{\rm mag}^{(N)}$ are $A_{\rm mag}$;
two physical ports cannot independently prescribe all $N$ modal
forces.

A complete conservative residual must now be specified as a physical
potential $\mathcal R_N(a^{(N)})$, measured in joules. The enriched
plant is
\begin{align}
 &\ddot a^{(N)}+\alpha\dot a^{(N)}+\Omega_N^2a^{(N)}\nonumber\\
 &\quad=\boldsymbol t_NE_0(\vartheta_N-\sin\vartheta_N)
       -\nabla_a\mathcal R_N\nonumber\\
 &\qquad+A_{\rm mag}^{(N)}\binom{F_{\rm app}}{\tau_{B,\rm app}}.
 \label{app:enriched-plant}
\end{align}
For a matched synthetic hierarchy, impose
$\mathcal R_N(D_sX,0)=E_0R_*(X)$ by restricting one declared full
potential to each modal subspace. This condition does not identify
that full potential from two-mode data. Supplement~\ref{sm:enriched}
gives explicit extensions and the low-mode projection.

Set $\vartheta_2=d^{\mathsf T}X$ and define the two additional terms
\begin{align}
 e_{\rm mag}^{(N)}
 &=d[\vartheta_N-\sin\vartheta_N\nonumber\\
 &\qquad-\vartheta_2+\sin\vartheta_2],\nonumber\\
 e_R^{(N)}
 &=\Lambda^{-1}\nabla_{a_{1:2}}\mathcal R_N(a^{(N)})\nonumber\\
 &\quad-\nabla R_*(X).
 \label{app:enriched-defects}
\end{align}
With the unchanged learned law, the exact low-mode equations become
\begin{align}
 X'&=\mathsf N Y,\nonumber\\
 Y'&=-\nabla(V_0+\delta R)(X)-\delta\gamma Y\nonumber\\
   &\quad+e_{\rm mag}^{(N)}-e_R^{(N)}+e_{\rm port}.
 \label{app:enriched-closed-loop}
\end{align}
The field conversion uses the actual full angle
$\vartheta=\vartheta_N$; replacing it by a two-mode estimate can also
create $e_{\rm port}$. Even exact port inversion does not eliminate the
high-mode angle or residual differences. A small two-mode training
loss alone therefore supplies neither a truncation bound nor the
Hamiltonian structure needed for the full plant's KAM analysis.

\subsection{A staged full-model validation protocol}
\label{app:staged-validation}
The next validation stages keep the fitted coefficients, low-mode
initial state, reference $h_0$ trajectory and port convention fixed.
The passive, damping-only, learned-compensation and truth-only groups
use the same plant, state information and actuator constraints at each
stage. The truth-only comparator retains its two-mode compensation
law; it is distinct from the stronger full-state feedback described
below. Zero initial high coordinates can still develop a response
because the same two ports act on every retained mode.
Table~\ref{app:validation-stages} separates the completed
numerical check from the proposed enrichments. Mesh convergence in
$16/32/64$ elements checks linear spatial discretization; changing the
number of retained modes checks the forced reduction and controller.

\begin{table*}[t]
\centering\small
\caption{Staged validation of the same fitted residual controller.
The higher-mode and independent finite-element stages are specified
protocols, not newly reported simulation results.}
\label{app:validation-stages}
\begin{tabular}{p{.15\textwidth}p{.45\textwidth}p{.30\textwidth}}
\toprule
Stage & Fixed object and controlled change & Status and required outputs\\
\midrule
Two modes & Fixed-kinetic residual class, one joint fit, common midpoint and magnetic interface & Completed: independent-reference errors, step convergence, reconstruction and sampled demands.\\
$4$ and $8$ modes & Same $32$-element basis, fixed learned law; declared full residual restricted consistently; first zero, then prescribed nonzero high-mode initial conditions & Planned: low-state error, high-mode response and the three defects in \eqref{app:enriched-closed-loop}.\\
All $64$ modes & Same finite-element mass, bias stiffness, core inertia and physical ports; full tip angle used in conversion & Planned: whole-centreline displacement and slope, input peaks, initialization and startup transients.\\
Independent FE model & Rebuild the same geometry, material assumptions, end mass/inertia and boundary loads in COMSOL or an independently specified solver & Planned: verify linear modes first, then compare driven responses and reduction errors under identical input definitions.\\
\bottomrule
\end{tabular}
\end{table*}

Each enriched run should record the full reconstructed centreline,
maximum sampled slope, minimum $|\cos\vartheta_N|$, state RMS error,
tip error, dominant and combination frequencies, force/field/gradient
peaks and saturation fraction. Compare time steps and spatial
extrema within each Hermite element; a plotted tip trace cannot bound
the rest of the beam. Sensor noise, damping mismatch and imposed
limits are introduced in separately declared tests with the same
constraints for every controller. A clipped command is applied to the
plant and included in $e_{\rm port}$, rather than retaining the
unclipped generalized force.

Supplement~\ref{sm:exact-lift} retains a conditional invariant-torus
construction for a stronger exact feedback using known damping and a
target acceleration. That theorem is useful when its full-state
hypotheses can be verified, but it does not replace the unchanged
learned controller in \eqref{app:enriched-closed-loop}. No new
$4/8/64$-mode forced trajectory, COMSOL comparison or finite-amplitude
torus certificate is claimed here.

\subsection{Whole-beam clearance from physical energy}
\label{app:clearance}
Let $E(t)=H_0(q,M_f\dot q)$ be the physical energy of
\eqref{mech:fe-plant}. For every clamped displacement,
\begin{equation}
 w(s)=\int_0^s(s-\xi)w_{ss}(\xi)\dd\xi.
 \label{app:clamped-identity}
\end{equation}
Cauchy--Schwarz and the nonnegative kinetic and bias energies give
\begin{equation}
 \|w\|_\infty^2\le\frac{2L^3E_{\rm bend}}{3EI}
                    \le\frac{2L^3E}{3EI},\qquad
 E_c=\frac{3EIc^2}{2L^3}.
 \label{app:energy-bound}
\end{equation}
Thus $E<E_c$ guarantees clearance of the entire centreline and its
radius tube in the declared straight corridor. The result holds for
arbitrary admissible displacement shapes, including all retained modes;
it is not restricted to a static tip-load shape. For this benchmark,
$E_c=10.0197171$ nJ and the selected cap
$E_{\max}=0.95E_c=9.5187312$ nJ.

An input-work governor can enforce this sufficient condition without
assuming the requested trajectory remains on a torus. Let
$u=(F,B_y)^{\mathsf T}$ and define the physical power coefficient
\begin{equation}
 a_P=\binom{\dot w(L)}{\mu\cos\vartheta\,\dot\vartheta},\qquad
 W(0)\ge E(0),\quad \dot W=a_P^{\mathsf T}u.
 \label{app:work-observer}
\end{equation}
By \eqref{mech:physical-power}, $E(t)\le W(t)$.
Given a request $u_{\rm req}$, choose the closest admissible input in
a scaled positive definite metric $R_u$:
\begin{align}
 u=\mathop{\arg\min}_{\widetilde u\in\mathcal U(q)}
       &\ \tfrac12(\widetilde u-u_{\rm req})^{\mathsf T}
                 R_u(\widetilde u-u_{\rm req}),\nonumber\\
 \text{subject to }&\quad
 a_P^{\mathsf T}\widetilde u\le\kappa(E_{\max}-W),
 \qquad \kappa>0.
 \label{app:work-governor}
\end{align}
Here $\mathcal U(q)$ is a nonempty closed convex feasible joint port set containing zero active
input; $R_u$ accounts for the different force and field units.
If $W(0)<E_{\max}$, zero input is feasible, and
$E_{\max}-W(t)\ge[E_{\max}-W(0)]e^{-\kappa t}>0$.
Together with \eqref{app:energy-bound}, this proves clearance for
solutions satisfying the power identity. Supplement~\ref{sm:work}
extends the argument to specified sensing and execution-error bounds.
For the enriched plant with residual $\mathcal R_N$, the derivative
of the nominal energy also contains
$-\dot a^{(N)\mathsf T}\nabla_a\mathcal R_N$. Its power must be
bounded and included in this ledger, or a separate lower bound of an
augmented energy by bending energy must be proved. The nominal
clearance theorem does not absorb that residual automatically. Activation of the governor can change the requested
Hamiltonian motion, so the clearance proof is distinct from torus
preservation.

The reported residual-compensation tests do not use this governor.
Their nominal-reference initial state has energy $12.6567$ nJ, already
above the $9.5187$ nJ cap.
Consequently, the observed displacement and slope maxima in
Section~\ref{sec:analysis} are numerical operating-domain diagnostics,
not consequences of this energy-cap guarantee. The governor supplies
a precise condition for a future full-state experiment initialized
inside the admissible energy set.
